\chapter{Evaluation}
\label{ch:evaluation}

\begin{guidance}[What this chapter does]
Critically assesses both the \emph{product} and the \emph{process}. This is
where the highest marks are earned: the Handbook asks whether the solution is
fit for purpose (\enquote{Is this suitable?}) and whether the work could have
been done better (\enquote{Could this have been done in a better way?}).
Compare the original theory and plan with what actually happened, and be
honest about limitations and flaws -- markers reward critical awareness, not
a project without problems. Typical length: 1,000--1,500 words.
\end{guidance}

\section{Evaluation of the product against the specification}
\label{sec:eval-product}

\begin{guidance}
Go through the requirements: met, partly met or not met, with the evidence
(test, figure, table) and the reasons for any discrepancy. A compact table
followed by discussion of the important points works well. How strong is the
evidence? What are the uncertainties?
\end{guidance}

\Cref{req:accuracy,req:interval} were met (\cref{tab:tests}).
\Cref{req:cost} was not met because \ldots

\section{Evaluation of the process}
\label{sec:eval-process}

\begin{guidance}
Compare the actual progress with the original plan (\cref{fig:gantt-summary})
and explain the deviations. Were the chosen methods, tools and techniques
appropriate? How well did the risk management work? What would you do
differently? Refer to the logbook as evidence.
\end{guidance}

\section{Limitations}
\label{sec:limitations}
